% Study guide for the methods of experiment 01 (replies x stance x lifetime, DerStandard)
% Build: pdflatex methods-study-guide.tex (twice)
\documentclass[11pt,a4paper]{article}
\usepackage[utf8]{inputenc}
\usepackage[T1]{fontenc}
\usepackage{lmodern}
\usepackage{amsmath,amssymb}
\usepackage[margin=2.4cm]{geometry}
\usepackage{booktabs,tabularx,array}
\usepackage{enumitem}
\usepackage{xcolor}
\usepackage{url}
\urlstyle{tt}
\usepackage[hidelinks]{hyperref}
\usepackage{microtype}

\setlist{itemsep=2pt,topsep=4pt}
\setlength{\parindent}{0pt}
\setlength{\parskip}{5pt}
\renewcommand{\arraystretch}{1.15}

\newcommand{\file}[1]{\path{#1}}
\newcommand{\inour}[1]{\par\textbf{In our code.} #1}
\newcommand{\pitfall}[1]{\par\textbf{Pitfall.} #1}
\newcommand{\readref}[1]{\par\textbf{Read.} #1}
\newcommand{\status}[1]{\textsc{\small #1}}
\newcolumntype{Y}{>{\raggedright\arraybackslash}X}
\newcommand{\OR}{\mathrm{OR}}
\newcommand{\E}{\mathbb{E}}

\title{\textbf{Study guide for the methods}\\[2pt]\large Experiment 01: who replies to whom, are they opposed, and how long do they stay?}
\author{State of 2026-11-02}
\date{}

\begin{document}
\maketitle
\tableofcontents

% ---------------------------------------------------------------
\section*{How to use this guide}
\addcontentsline{toc}{section}{How to use this guide}
The goal is that you can \emph{explain, defend and criticise} each method of the experiment without the code in front of you. For each method the guide gives the same blocks: the \textbf{question} it answers in this project, the \textbf{idea} in plain words, a \textbf{worked example with small numbers} you can redo by hand, \textbf{where it is in our code}, the usual \textbf{pitfalls} (including the ones found in the code review), and what to \textbf{read}.

This guide explains the methods \emph{as needed to use them correctly}; it is not a replacement for the textbooks. Sources are in the reading list at the end. Section~\ref{sec:selftest} has questions with model answers; try to answer them aloud before reading the answers.

\textbf{Numbers.} Every result quoted here comes from \file{notes/experiments/01-reply-opposition-lifetime/RESULTS.md} (section 8 for the positive-vote and long-lived follow-ups) or the review in \file{CODE-REVIEW.md}. The toy examples are mine and are computed exactly; they are not data.

\textbf{Suggested path (about 10 hours).}\par\smallskip
\begin{tabularx}{\textwidth}{@{}l Y r@{}}
\toprule
Block & Contents & Time \\
\midrule
1 & Sections~\ref{sec:map}--\ref{sec:vocab}: the map and the vocabulary & 1\,h \\
2 & Sections~\ref{sec:or}--\ref{sec:boot}: H1, H2 (odds ratios, stratification, null, Beta, bootstrap) & 3\,h \\
3 & Section~\ref{sec:signed}: H3 (signed networks, factions) & 1.5\,h \\
4 & Sections~\ref{sec:surv}--\ref{sec:h4}: lifetime and H4 (censoring, Kaplan--Meier, landmark, Cox) & 3\,h \\
5 & Sections~\ref{sec:robust}--\ref{sec:claims}: robustness, follow-ups, what can be claimed & 1\,h \\
6 & Section~\ref{sec:selftest}: self-test & 0.5\,h \\
\bottomrule
\end{tabularx}

% ---------------------------------------------------------------
\section{The map}\label{sec:map}
\subsection{From question to measure}
There is \emph{no comment text}, so ``position'' is inferred from \textbf{votes}. Everything below is therefore \emph{vote-based disapproval or approval}, not ideology.

\begin{tabularx}{\textwidth}{@{}l Y Y@{}}
\toprule
Hyp. & Compared groups & Statistic \\
\midrule
H1 & On the same comment: voters who \emph{replied} to it vs other voters & Mantel--Haenszel odds ratio of a downvote \\
H2 & A's overall votes toward the person A replied to (B) vs toward a random co-present author (the null) & Share of negative relations; mean posterior approval; paired bootstrap \\
H3 & Replies where A and B belong to the two FAULTANA factions: share crossing factions vs random mixing and vs the null & Cross-faction share; downvote rate by relation \\
H4 & New commenters, by their first-30-day exposure to downvotes and hostile replies & Hazard ratio of leaving (Cox, landmark at day 30) \\
\bottomrule
\end{tabularx}

\subsection{Data flow in one picture (words)}
Raw TSV $\to$ typed Parquet (\file{00}) $\to$ reply edges A$\to$B (\file{01}) $\to$ null target per reply (\file{02}) $\to$ vote totals per pair and per edge (\file{03}) $\to$ tests H1 (\file{04}), H2/H3 (\file{05}), lifetime and H4 (\file{06}), robustness (\file{07}), positive votes (\file{08}), long-lived users (\file{09}). The shared H4 feature builder is \file{h4.py}.

\subsection{Three levels of ``opposition''}
\emph{Comment} (does A disapprove of \emph{this} comment?), \emph{pair} (is A \emph{generally} against B?), \emph{faction} (are A and B on opposite sides of the platform's fault line?). They answer different questions and can disagree; if they all point the same way, the result is not specific to one level.

% ---------------------------------------------------------------
\section{Vocabulary you must be fluent in}\label{sec:vocab}
\begin{description}[style=nextline,leftmargin=0pt]
\item[Probability, odds, odds ratio.] If a share is $p$, the \emph{odds} are $p/(1-p)$. The \emph{odds ratio} compares two groups: $\OR=\dfrac{p_1/(1-p_1)}{p_0/(1-p_0)}$. $\OR=1$: no difference; $>1$: group 1 has higher odds.
\textbf{Worked example (H1 shares).} Repliers downvote 53.3\% and other voters 30.1\%. Odds: $0.533/0.467=1.14$ and $0.301/0.699=0.43$. Crude $\OR=1.14/0.43=2.65$. The \emph{relative risk} is $0.533/0.301=1.77$. So ``4.4 times more likely'' would be wrong twice: the 4.43 is an odds ratio (not a risk ratio), and it is the \emph{stratified} value (Section~\ref{sec:or}).
\item[Rate and hazard.] The hazard is the instantaneous risk of the event (here, leaving) among those still present. A \emph{hazard ratio} (HR) compares hazards; HR $=1.11$ per SD means: at any moment, a user one SD higher in the exposure has an 11\% higher hazard of leaving than an otherwise comparable user.
\item[Confidence interval (CI).] A range of values compatible with the data under the model's assumptions. It describes \emph{sampling uncertainty}, not bias.
\item[Effect size vs $p$-value.] With millions of rows every difference is ``significant''; the information is in the size of the effect and its robustness. \readref{\cite{wasserstein}, \cite{lin}.}
\item[Association vs causation.] A difference observed in data, with no control over what caused it. Our results are associations (votes carry no timestamp; no randomization). \readref{\cite{hernan}.}
\end{description}

% ---------------------------------------------------------------
\section{H1: odds ratio and Mantel--Haenszel}\label{sec:or}
\subsection*{Question}
Do people who reply to a comment downvote it more than other people who voted on the \emph{same} comment?

\subsection*{Idea}
For one comment, build the $2\times2$ table
\[
\begin{array}{c|cc}
 & \text{downvote} & \text{upvote}\\\hline
\text{voter replied} & a & b\\
\text{voter did not reply} & c & d
\end{array}
\qquad \OR=\frac{ad}{bc}.
\]
Different comments have different baselines (a provocative comment gets many downvotes from everyone). Pooling all votes into one table would mix those baselines. \textbf{Stratifying by comment} compares like with like, and the \textbf{Mantel--Haenszel (MH) estimator} combines the per-comment tables into one common odds ratio:
\[
\widehat{\OR}_{MH}=\frac{\sum_k R_k}{\sum_k S_k},\quad R_k=\frac{a_kd_k}{n_k},\ S_k=\frac{b_kc_k}{n_k},\ n_k=a_k+b_k+c_k+d_k .
\]
Each comment contributes with weight proportional to its information; comments where one cell group is empty contribute nothing, so we keep only comments with at least one replier-voter and one other voter.

\subsection*{Worked example (two comments)}
\begin{tabular}{@{}lcccc@{\hspace{1.5em}}ccc@{}}
\toprule
Comment & $a$ & $b$ & $c$ & $d$ & $\OR_k$ & $R_k$ & $S_k$\\
\midrule
X (controversial) & 6 & 2 & 24 & 16 & 2.00 & 2.00 & 1.00\\
Y (calm) & 1 & 9 & 4 & 36 & 1.00 & 0.72 & 0.72\\
\bottomrule
\end{tabular}

Pooling the two tables gives $a=7,b=11,c=28,d=52$ and a crude $\OR=364/308=1.18$. MH gives $(2.00+0.72)/(1.00+0.72)=1.58$. The two numbers differ because the crude table lets comment X, which has more downvotes overall, dominate the pooled counts; MH compares repliers with non-repliers \emph{inside} each comment. Redo this by hand until you can do it without looking.

\subsection*{Uncertainty}
The CI uses the Robins--Breslow--Greenland variance (a standard variance for the MH log odds ratio; \cite{rbg}). It assumes the votes are independent. A second interval is a bootstrap that resamples \emph{comments} with Poisson(1) weights. Neither accounts for the fact that the \emph{same voter} casts many votes, so both are too narrow (review item M-1).

\subsection*{Results and how to read them}
Pooled $\OR_{MH}=4.43$ over 8.1M comments; downvote shares 53.3\% (repliers) vs 30.1\% (others). By thread depth of the replied comment: OR 3.5 (top level) $\to$ 9.0 (depth $\ge5$). The OR does not say ``4.4 times''; quote the shares beside it.

\subsection*{Robustness check run in the review}
Repliers may simply be different voters (heavier, more negative). I re-did H1 with \textbf{voters} as strata: each voter's votes on comments they replied to vs their votes on other comments that had received at least one reply. $\OR_{MH}=3.49$; 89\% of voters with $\ge10$ votes on both sides downvote more on the comments they reply to. So voter heterogeneity explains part of 4.43 (about 16\% on the log scale) but not the effect.

\inour{\file{04_h1_comment_level.py} (cells and MH in \texttt{mh()}, strata query, bootstrap). Positive-vote mirror: \file{08_positive_votes.py}; since up and down are the two columns of the same table, $\OR_{up}=1/\OR_{down}$ exactly (0.23 = 1/4.43), so it is not independent evidence.}
\pitfall{We do not know whether the downvote came before or after the reply. Treat as association. Also: a comment with no replier-voter and one other voter contributes nothing, so MH describes the comments that attracted replies.}
\readref{\cite{mh}, \cite{rbg}; \cite{agresti} (stratified $2\times2$ tables).}

% ---------------------------------------------------------------
\section{H2: user pairs, the null model and Beta edges}\label{sec:null}
\subsection*{Question}
Does A reply to people toward whom A is, in general, more negative than toward other people who were present?

\subsection{Reply edges}
A \textbf{reply edge} A$\to$B is a comment by A whose parent comment is written by B. We drop self-replies, replies involving deleted users (author id missing), and records flagged as inconsistent. Only the thread structure certifies a reply; a free-text mention of B is not counted.

\subsection{The co-present null}
\textbf{Idea.} To say that 31\% negative relations is ``high'', we need what we would see if A had chosen differently. For each reply, the \textbf{null target} is the author of a uniformly random \emph{earlier} comment in the \emph{same article} (not A herself). Same article, same moment, same exposure; only the \emph{choice of target} is randomised. Then every H2 statistic is a \textbf{paired comparison}: for the same replier A, real target B vs null target $B'$.

\textbf{What it controls:} article, time and the fact that the null comment existed. \textbf{What it does not control} (review M-3): the real target is chosen among comments A actually \emph{saw}, which concentrate in A's own thread and in more visible comments. A second null at \emph{thread} level would remove that. The correct wording is therefore ``more negative than toward a random co-present commenter in the article''.

\subsection{Dyad votes and the Beta posterior}
For a directed pair (A,B), count A's positive and negative votes on \emph{all} B's comments in the whole dataset: $\text{pos},\text{neg}$. Treat the unknown probability $p$ that A votes positively on B as a Beta$(\alpha,\beta)$ variable, updated by the counts:
\[
\E[p]=\frac{\text{pos}+\alpha}{\text{pos}+\text{neg}+\alpha+\beta},\qquad
\mathrm{Var}[p]=\frac{\E[p](1-\E[p])}{\text{pos}+\text{neg}+\alpha+\beta+1}.
\]
We inherit from FAULTANA the prior $(\alpha,\beta)=(1,2)$ (mildly favouring negative, because negative votes are rarer) and the rule: edge \textbf{negative} if $\E[p]<0.4$, \textbf{positive} if $>0.6$ \cite{faultana}.

\textbf{Worked example.}\par\smallskip
\begin{tabular}{@{}rrrrp{7.2cm}@{}}
\toprule
pos & neg & $\E[p]$ & $\mathrm{Var}$ & reading\\
\midrule
1 & 3 & 0.286 & 0.0255 & below 0.4, but only 4 votes (undefined under ``$\ge5$ votes'')\\
2 & 8 & 0.231 & 0.0127 & negative edge\\
50 & 50 & 0.495 & 0.0024 & no edge (inside 0.4--0.6)\\
400 & 100 & 0.797 & 0.00032 & positive edge\\
\bottomrule
\end{tabular}

The prior matters with few votes (one downvote moves $\E[p]$ to $0.25$). That is why we also require a minimum number of votes (5 in the headline) and vary it (1, 5, 20) in the robustness analysis. The FAULTANA ``strict'' rule additionally needs $\mathrm{Var}<10^{-4}$, which at $\E[p]\approx0.5$ requires thousands of votes; it selects heavy pairs.

\subsection{What is compared}
Among replies with a null target, per replier we compute: the share of \textbf{negative edges among defined edges} (edges with $\ge5$ votes and $\E[p]$ outside $[0.4,0.6]$), the mean $\E[p]$, and the share with any vote history. Headline (A's vote on the replied comment excluded, so H2 does not re-count H1): negative edges 31.4\% (real) vs 18.1\% (null), $+13.3$ pp [12.6, 14.0].

\inour{\file{02_null_pairs.py} (null draw: rank $k$ of the reply in the article, $k_{null}=1+\lfloor U(k-1)\rfloor$ with $U$ from a hash), \file{03_dyads.py} (vote totals), \file{05_h2_h3_dyad_faction.py} (\texttt{stance\_cols}). The ``excl'' variant subtracts A's vote on the replied (or null) comment.}
\pitfall{(i) ``Defined edges'' are a selected subset and real pairs have more shared history than null pairs (any votes: 59.4\% vs 55.6\%), so part of the gap is a coverage effect (review M-2). (ii) Dyad votes include votes after the reply, possibly caused by it. (iii) The depth breakdown used pooled shares, which the notes call misleading (review M-4).}
\readref{\cite{faultana} (Methods: network creation), \cite{bda} (beta--binomial).}

% ---------------------------------------------------------------
\section{Bootstrap clustered by replier}\label{sec:boot}
\textbf{Question.} How uncertain is ``real $-$ null'' when one replier contributes many correlated edges?

\textbf{Idea.} Resample \emph{users}, not rows. Each replier gets a random integer weight $w_A\sim\mathrm{Poisson}(1)$ (the Poisson bootstrap, a fast approximation to resampling users with replacement); recompute the statistic as a ratio of weighted sums; repeat 200 times; use the 2.5th--97.5th percentiles of the \emph{difference} real$-$null. Because real and null share the same weights, the interval is for the paired difference.

\textbf{Worked example.} Three repliers with (negative edges, defined edges) $=(2,10),(5,10),(1,10)$ in the real condition give $8/30=26.7\%$. Draw weights $(2,0,1)$: $(2\cdot2+0+1)/(2\cdot10+0+10)=5/30=16.7\%$. Doing this many times shows how much the statistic moves when users are reweighted.

\inour{\file{05_h2_h3_dyad_faction.py}: \texttt{per\_a}, \texttt{summarise(w)}, 200 replicates. H1 uses weights over comments, not voters (so it is not user-clustered; review M-1).}
\readref{\cite{efron}.}

% ---------------------------------------------------------------
\section{H3: signed networks, balance and factions}\label{sec:signed}
\subsection*{Question}
Do replies cross the platform's fault line less (or more) than chance, and are cross-faction replies more hostile?

\subsection*{Idea 1: balance (Heider; Cartwright and Harary)}
A signed relation is $+$ (like) or $-$ (dislike). Heider's triad rule: ``a friend of my friend is my friend; an enemy of my enemy is my friend''; a triangle is balanced if the product of its signs is positive (all $+$, or two $-$ and one $+$) \cite{heider}. Cartwright and Harary generalised it to graphs: a signed graph is balanced iff every cycle has a positive sign, and (\emph{structure theorem}) iff the nodes can be split into \textbf{two camps}, positive edges inside and negative edges between \cite{cartwright}. They also define a \emph{degree of balance} (share of positive cycles) and stress that \emph{no tie is not a negative tie}.

\subsection*{Idea 2: FAULTANA}
Real networks are not perfectly balanced. FAULTANA finds the partition that minimises \emph{frustrated edges} (negative inside a group, positive between groups), then measures Alignment and Antagonism. On DerStandard it finds $k^*=2$ groups for 14.8k stable users; our \file{faction\_users} file gives each user's group (about 62\% in the larger one) \cite{faultana}.

\subsection*{What H3 computes}
Among replies where both users are labelled (11.3M): cross-faction share observed 26.0\%, null 23.2\%, random mixing by faction size 46.8\% ($2p(1-p)$ with $p=0.626$). Reading: users mostly share discussions with their own faction (exposure is segregated), and within that exposure choose out-group targets slightly more than chance. Cross-faction replies are hostile: the replier downvoted the parent in 64.6\% of replies vs 37.0\% within a faction.

\inour{\file{05_h2_h3_dyad_faction.py} (\texttt{h3\_factions}). The projection of labels to all users (AUC 0.63, in-sample) is weak and is not used in the conclusions.}
\pitfall{\textbf{Circularity}: the factions were derived from the same vote network, so ``cross-faction replies are downvoted more'' is partly built in. \textbf{Survivors}: only users who voted every year 2014--2021 have a label. \textbf{Static}: the labels do not change over time.}
\readref{\cite{heider}, \cite{cartwright}, \cite{faultana}; \cite{easley} (chapter on positive and negative relationships).}

% ---------------------------------------------------------------
\section{Lifetime and survival analysis}\label{sec:surv}
\subsection*{Question}
How long do users stay, and how does that differ between cohorts?

\subsection*{Definitions}
\textbf{Lifetime} $=$ last activity $-$ first activity (comments and votes). The data end on 2022-12-31, so a user whose last activity is after 2022-07-01 may simply still be active: \textbf{right-censored} (the 6-month buffer is a convention for ``churned''). A user whose first activity is before 2013-07-01 may have started before the data: \textbf{left-censored} (excluded from cohort analyses). \emph{Cohort} $=$ entry year.

\subsection*{Kaplan--Meier}
$\hat S(t)$ is the estimated share still active after time $t$, counting censored users as at risk until their censoring time. At each event time $t_i$ with $n_i$ users at risk and $d_i$ events: $\hat S(t)=\prod_{t_i\le t}(1-d_i/n_i)$.

\textbf{Worked example.} Six users with (months, left?): $(1,\text{yes}),(2,\text{yes}),(2,\text{no}),(4,\text{yes}),(5,\text{no}),(6,\text{no})$ where ``no'' means still active (censored).\par\smallskip
\begin{tabular}{@{}rrrr@{}}
\toprule
$t$ & at risk & events & $\hat S$\\
\midrule
1 & 6 & 1 & $5/6=0.833$\\
2 & 5 & 1 & $0.833\cdot4/5=0.667$\\
4 & 3 & 1 & $0.667\cdot2/3=0.444$\\
\bottomrule
\end{tabular}

The censored user at $t=2$ leaves the risk set after $t=2$ without being counted as an event. Dropping censored users or treating them as events would bias the curve.

\subsection*{Results}
27\% of users are active for less than a day, 38\% for less than 30 days; Kaplan--Meier median lifetime of commenters falls by cohort: 1.36 y (2014), 0.81 (2016), 0.41 (2018), 0.39 (2020).

\inour{\file{06_lifetime.py} (descriptives, CCDF, Kaplan--Meier by cohort).}
\pitfall{\textbf{Deleted users are not in the users file} (their ids are missing in comments and votes), so exits by account deletion or ban are not observed; lifetimes describe users who kept their account (review H-1). Votes have no timestamp, so ``last activity'' may be a vote on an old comment.}
\readref{\cite{km}, \cite{kleinbaum}.}

% ---------------------------------------------------------------
\section{H4: landmark, Cox model and shrinkage}\label{sec:h4}
\subsection*{Question}
Do new commenters who receive more downvotes and hostile replies in their first 30 days leave sooner?

\subsection{Why a landmark: immortal-time bias}
Suppose we measured exposure over each user's \emph{whole} stay and related it to the length of the stay. Users who stay longer have more time to accumulate replies and downvotes, so ``more exposure'' trivially goes with ``longer stay''; the time before exposure is complete is ``immortal'' (the user could not have left and still been counted in that exposure group). The \textbf{landmark} design fixes a calendar point (day 30 after the first comment), measures exposures \emph{before} it, and counts survival \emph{after} it, only among users still active at the landmark \cite{suissa,anderson}.

\textbf{Worked example.} User P writes 2 comments on day 1 and leaves on day 10; user Q writes 2 comments on day 1 and stays 200 days. At the day-30 landmark, P is not in the analysis (already gone), Q is; P's early exit is reported separately (gone before the landmark: a descriptive share). The comparison then uses only users who \emph{could} be followed from day 30.

\subsection{Features and shrinkage}
For each user, in the first 30 days: downvote share received, hostile replies received (replier downvoted the user's comment), hostile replies written, whether any reply was received, and log of the number of comments. Shares from few events are noisy, so they are \textbf{shrunk} toward the population rate $p_0$ with pseudo-count $M$: $\tilde p=(k+Mp_0)/(n+M)$.

\textbf{Worked example} ($p_0=0.17$, $M=2$):\par\smallskip
\begin{tabular}{@{}rrr@{}}
\toprule
hostile $k$ & replies $n$ & $\tilde p$\\
\midrule
0 & 0 & 0.170 (the population rate)\\
1 & 1 & 0.447\\
0 & 1 & 0.113\\
5 & 10 & 0.445\\
50 & 100 & 0.494\\
\bottomrule
\end{tabular}

A first version shrank toward 0.5 (a Beta(1,1) prior), which put users with \emph{no} data at 0.5, above the real base rate; ``no data'' looked hostile and the sign of a coefficient changed. Shrinking toward the population rate fixed it.

\subsection{The Cox proportional-hazards model}
The hazard of leaving at time $t$ for a user with covariates $x$ is $h(t\mid x)=h_0(t)\exp(\beta^\top x)$. The baseline $h_0(t)$ is left unspecified; the model estimates only the \textbf{hazard ratios} $\exp(\beta_j)$. We standardise continuous covariates, so $\exp(\beta_j)$ is the HR per one standard deviation; for the binary ``got any reply'' it is the HR yes vs no. Two users 2 SD apart differ by $\mathrm{HR}^2$ (for $1.11$: $1.23$).

\textbf{Stratification by cohort} gives each entry cohort its own baseline hazard, because newer cohorts leave faster; the coefficients are common to all cohorts. \textbf{Proportional hazards} means the HR does not change over time; it is checked with Schoenfeld residuals \cite{schoenfeld,grambsch}. Our $\mathrm{PH}$ test used a 20k subsample (review L-2).

\textbf{Results.} 98,670 new commenters (59\% leave before 2022-07). Per SD: downvote share received $\mathrm{HR}=1.11$ [1.10, 1.12]; hostile replies received $1.06$; hostile replies written $1.04$; got any reply $0.95$; log early comments $1.23$ (very active new accounts leave fast; probably bursts or banned accounts). Same sign in all 18 specifications.

\inour{\file{h4.py} (population, features, shrinkage, landmark) and \file{06_lifetime.py} (Cox with \texttt{CoxPHFitter(penalizer=0.001), strata=cohort}). The small penalizer is a tiny ridge shrinkage.}
\pitfall{(i) HR $>1$ means leaving \emph{sooner}. (ii) Votes carry no timestamp, so ``early'' downvotes may have been cast after day 30 (review H-3). (iii) The hostile share is relative to \emph{all} replies, but only 25.5\% of replies have an observed vote, so it mixes hostility with voting propensity (review M-5). (iv) Observational: reverse causation is possible (users already disengaging may provoke or attract hostility).}
\readref{\cite{cox}, \cite{kleinbaum}, \cite{suissa}, \cite{anderson}, \cite{lifelines}.}

% ---------------------------------------------------------------
\section{Robustness: multiverse and specification curves}\label{sec:robust}
\textbf{Idea.} Several choices are arbitrary but defensible: the prior $(1,2)$ vs $(1,1)$, the thresholds $0.6/0.4$ vs $0.7/0.3$, the minimum votes (1, 5, 20), whether to exclude the specific vote, the churn buffer (3, 6, 12 months), the early window (14, 30, 60 days), the shrinkage $M$. Instead of reporting one combination, run \emph{all} reasonable ones and show the spread \cite{steegen}. A specification curve additionally plots the sorted estimates with a ``dashboard'' of choices and adds a \emph{joint test} by resampling under the null \cite{simonsohn}.

\textbf{What we did.} H2: 24 specifications, gap between $+11.7$ and $+20.6$ pp, positive in all. H4: 18 specifications, HR for downvote share between 1.10 and 1.12. Positive votes: 24 specifications, all negative.

\pitfall{Our robustness is \emph{descriptive}; we ran no joint inference, and the H4 grid does not vary the model. Say ``same sign and similar size across 24 and 18 specifications'', not ``inferentially robust''. A specification should be a valid test of the question, not just a different number.}
\readref{\cite{steegen}, \cite{simonsohn} (both are in \file{notes/papers/}).}

% ---------------------------------------------------------------
\section{Follow-ups: positive votes and the longest-living users}
\subsection{Positive votes}
The positive-vote analysis mirrors the negative one with the same data, prior, thresholds, null and specifications. At the comment level it adds no information (the OR is the inverse). At pair level, replied-to users have a lower positive-edge share than the null (68.6\% vs 81.9\%, $-13.3$ pp), consistent with H2. In H4, the \emph{upvote share received} (HR 0.87 per SD) is the complement of the downvote share, so it carries the same information; \emph{friendly replies} have HR $\approx1.00$ once hostility is controlled. Reading: what matters is the \emph{absence of hostility}, not the presence of friendliness. This is an inference from the model; a ``friendly reply'' also requires an observed upvote.

\subsection{Longest-living users}
Population: the H4 landmark users. Groups: top 10\% of lifetime within cohort (preferred), top 10\% overall, $>2$ y, $>3$ y. Results: long-lived users received fewer downvotes (20.4\% vs 30.0\%) and wrote fewer hostile replies (12.2\% vs 17.1\%) in the first 30 days, with the same early activity and the same chance of getting a reply; per SD of downvote share received, the odds of being long-lived are 0.87.

\pitfall{\textbf{Redefine before using in the thesis (review H-2).} Lifetime is right-censored, so the top 10\% within cohort is entirely made of users still active at the end, ordered by \emph{when they last happened to be active}; and for fixed thresholds the comparison group includes cohorts that cannot reach the threshold. Compare within cohort, restrict fixed thresholds to eligible cohorts, and add a volume check.}

% ---------------------------------------------------------------
\section{What can be claimed}\label{sec:claims}
\begin{tabularx}{\textwidth}{@{}Y Y@{}}
\toprule
Can claim & Cannot claim \\
\midrule
On the same comment, repliers disapprove (by vote) of it much more than other voters, with and without controlling for who the voter is (OR 4.43 by comment, 3.49 by voter). & That disagreement \emph{causes} replying: votes are untimed. \\
Replied-to users are more negative (and less positive) than random co-present authors in the article. & That people ``choose enemies'': thread-level exposure and visibility are not controlled (null is article-level). \\
Early downvotes and hostile replies are associated with a higher hazard of leaving, with the same sign across 18 specifications. & That hostility \emph{causes} leaving, or that veterans ``gatekeep''; exits by deletion or ban are not observed. \\
Factions rarely meet; when they do, replies are mostly downvoted. & That this is independent evidence: factions come from the same votes. \\
\bottomrule
\end{tabularx}

\subsection*{Frequent confounders to have ready}
Voter heterogeneity (checked), visibility and popularity of the parent comment, thread-level exposure, activity volume (heavy users do everything more), cohort/calendar effects (stratified in H4), deleted accounts (unobserved), the 2022 volume drop (unexplained), vote timing (unobservable).

% ---------------------------------------------------------------
\section{Self-test with model answers}\label{sec:selftest}
Answer aloud first.
\begin{enumerate}
\item \textbf{Why can't ``position'' be measured from text here, and what replaces it?} The dataset has no comment text, only votes; a downvote is read as disapproval of that comment, aggregated at comment, pair and faction level.
\item \textbf{What is a reply edge and what is lost with that definition?} A comment whose parent is B's comment. Mentions without a parent link, and replies involving deleted users or inconsistent records, are lost.
\item \textbf{Why compare with a co-present null rather than random users?} Random users were not there; the null keeps article, time and exposure and changes only the choice of target. It does not keep the thread or comment visibility.
\item \textbf{Why is $\OR=4.43$ not ``4.4 times more likely''?} It is a ratio of odds; the shares are 53.3\% vs 30.1\% (risk ratio 1.77, crude odds ratio 2.65); 4.43 is the stratified value.
\item \textbf{What does stratifying by comment do?} It compares repliers and others on the same comment so that controversial comments do not inflate the effect; MH pools the per-comment ORs.
\item \textbf{What did the voter-stratified check show?} OR 3.49: voter heterogeneity explains part of the 4.43 but 89\% of voters still downvote more on comments they reply to.
\item \textbf{Why exclude A's vote on the replied comment in H2?} Otherwise H2 re-counts H1: the pair stance would include the very vote that H1 studies.
\item \textbf{What does the Beta prior do and why a minimum of votes?} It stabilises small samples (one downvote gives $\E[p]=0.25$); the minimum number of votes avoids calling ``negative'' a pair with a single vote.
\item \textbf{Why a Poisson bootstrap by replier?} Edges of the same replier are correlated; weighting users resamples users; paired weights give a paired CI.
\item \textbf{What is balance, and how does it relate to factions?} A signed graph is balanced iff it splits into two camps with positive inside and negative between; factions are the two camps that best fit the data (fewest frustrated edges).
\item \textbf{Why is H3b partly circular?} The faction partition was built from the same votes whose sign we then relate to cross-faction replies.
\item \textbf{What is right-censoring and what does it do to the top 10\%?} A user still active at the end may stay longer; the observed lifetime is a lower bound; the top group is mostly such users, ordered by last-activity date.
\item \textbf{What is immortal-time bias and how does the landmark avoid it?} Counting exposure over a period in which the user must have survived makes survivors look exposed; the landmark measures exposure before day 30 and follows only users alive then.
\item \textbf{What does HR $=1.11$ per SD mean, and what does it not mean?} At any time, a user one SD higher in downvote share received has an 11\% higher hazard of leaving than an otherwise comparable user; not a causal effect or a probability.
\item \textbf{Why stratify the Cox model by cohort?} Cohorts differ in baseline hazard; stratification gives each its own baseline instead of assuming one.
\item \textbf{What does the multiverse show, and what does it not?} The sign and size are stable across reasonable choices; it is not a formal test and the grid does not include model choices.
\item \textbf{Which limitations would you volunteer first?} Votes untimed (no causal order, early exposure may include later votes); deleted users unobserved; null at article level; factions circular; observational design.
\end{enumerate}

% ---------------------------------------------------------------
\section{Reading path}
Read in this order; the first item is the one that corrects most misunderstandings.
\begin{enumerate}
\item Hernán and Robins, \emph{Causal Inference: What If}, chapter 1 \cite{hernan}: association vs causation.
\item The dataset paper \cite{dataset} and FAULTANA \cite{faultana} (summaries in \file{notes/papers/PAPERS-SUMMARY.md}).
\item Agresti, stratified tables and odds ratios \cite{agresti}; Mantel and Haenszel \cite{mh}.
\item Kleinbaum and Klein, censoring, Kaplan--Meier, Cox \cite{kleinbaum}; then Suissa on immortal-time bias \cite{suissa}.
\item Heider and Cartwright--Harary \cite{heider,cartwright}; Easley and Kleinberg, signed-networks chapter \cite{easley}.
\item Steegen and Simonsohn \cite{steegen,simonsohn}; Wasserstein and Lazar on $p$-values \cite{wasserstein}.
\end{enumerate}
\textbf{Caveat.} The bibliographic details were written from memory and the existence of most sources was checked by web search on 2026-11-02 (see \file{PAPERS-SUMMARY.md}, section A); Efron 1979, Schoenfeld 1982 and a few textbooks were not verified. Check DOIs and pages before they enter the written thesis.

\begin{thebibliography}{99}\small
\bibitem{dataset} Fraxanet, G\'omez, Kaltenbrunner, Pellert (2026). A Decade of News Forum Interactions: Threaded Conversations, Signed Votes, and Topical Tags. \emph{Scientific Data}, doi:10.1038/s41597-026-07267-3.
\bibitem{faultana} Fraxanet, Pellert, Schweighofer, G\'omez, Garcia (2024). Unpacking polarization: Antagonism and alignment in signed networks of online interaction. \emph{PNAS Nexus} 3, pgae276.
\bibitem{heider} Heider, F. (1946). Attitudes and cognitive organization. \emph{J. Psychology} 21, 107--112.
\bibitem{cartwright} Cartwright, D., Harary, F. (1956). Structural balance: a generalization of Heider's theory. \emph{Psychological Review} 63(5), 277--293.
\bibitem{easley} Easley, D., Kleinberg, J. (2010). \emph{Networks, Crowds, and Markets}. Cambridge University Press.
\bibitem{mh} Mantel, N., Haenszel, W. (1959). Statistical aspects of the analysis of data from retrospective studies of disease. \emph{J. National Cancer Institute} 22, 719--748.
\bibitem{rbg} Robins, J., Breslow, N., Greenland, S. (1986). Estimators of the Mantel--Haenszel variance consistent in both sparse data and large-strata limiting models. \emph{Biometrics} 42, 311--323.
\bibitem{agresti} Agresti, A. (2013). \emph{Categorical Data Analysis}, 3rd ed. Wiley.
\bibitem{efron} Efron, B. (1979). Bootstrap methods: another look at the jackknife. \emph{Annals of Statistics} 7, 1--26.
\bibitem{bda} Gelman, A., et al. (2013). \emph{Bayesian Data Analysis}, 3rd ed. CRC Press.
\bibitem{km} Kaplan, E. L., Meier, P. (1958). Nonparametric estimation from incomplete observations. \emph{J. American Statistical Association} 53, 457--481.
\bibitem{cox} Cox, D. R. (1972). Regression models and life-tables. \emph{J. Royal Statistical Society B} 34, 187--220.
\bibitem{kleinbaum} Kleinbaum, D., Klein, M. \emph{Survival Analysis: A Self-Learning Text}. Springer.
\bibitem{schoenfeld} Schoenfeld, D. (1982). Partial residuals for the proportional hazards regression model. \emph{Biometrika} 69, 239--241.
\bibitem{grambsch} Grambsch, P., Therneau, T. (1994). Proportional hazards tests and diagnostics based on weighted residuals. \emph{Biometrika} 81, 515--526.
\bibitem{anderson} Anderson, J., Cain, K., Gelber, R. (1983). Analysis of survival by tumor response. \emph{J. Clinical Oncology} 1, 710--719.
\bibitem{suissa} Suissa, S. (2008). Immortal time bias in pharmaco-epidemiology. \emph{American J. Epidemiology} 167, 492--499.
\bibitem{steegen} Steegen, S., Tuerlinckx, F., Gelman, A., Vanpaemel, W. (2016). Increasing transparency through a multiverse analysis. \emph{Perspectives on Psychological Science} 11, 702--712.
\bibitem{simonsohn} Simonsohn, U., Simmons, J., Nelson, L. (2020). Specification curve analysis. \emph{Nature Human Behaviour} 4, 1208--1214.
\bibitem{wasserstein} Wasserstein, R., Lazar, N. (2016). The ASA statement on p-values: context, process, and purpose. \emph{The American Statistician} 70, 129--133.
\bibitem{lin} Lin, M., Lucas, H., Shmueli, G. (2013). Too big to fail: large samples and the p-value problem. \emph{Information Systems Research} 24, 906--917.
\bibitem{hernan} Hern\'an, M., Robins, J. (2020). \emph{Causal Inference: What If}. Chapman \& Hall/CRC.
\bibitem{lifelines} Davidson-Pilon, C. (2019). lifelines: survival analysis in Python. \emph{J. Open Source Software} 4(40), 1317.
\end{thebibliography}

\end{document}
